\section*{Bab \ref{ch:b3:submanifolds} --- Submanifold
$\R^n$}

\begin{solution}{exo:b3:submanifolds:1}
(a) Masing-masing berupa $F^{-1}(0)$ bagi sebuah submersi: $\norm x^2 - 1$ pada
$\R^n\setminus\{0\}$ (dengan gradien $2x \neq 0$): jadi berdimensi $n-1$;
$x^2 + y^2 - z^2 - 1$ (dengan gradien $(2x, 2y, -2z) \neq 0$ pada
hiperboloidnya, yang $x^2 + y^2 = 1 + z^2 > 0$ padanya): jadi berdimensi $2$;
sedangkan fungsi torusnya $G = (\rho - R)^2 + z^2 - r^2$ dengan $\rho =
\sqrt{x^2+y^2}$ bersifat $\mathcal C^\infty$ di dekat torusnya (sebab $\rho
\geq R - r > 0$ di sana) dengan $\nabla G \neq 0$ (sebab
komponen-$z$-nya adalah $2z$, dan di tempat $z = 0$ komponen radialnya
adalah $2(\rho - R)\ne0$ karena $\abs{\rho - R} = r$): jadi berdimensi
$2$.

(b) Untuk $\varepsilon$ yang kecil,
$(C\setminus\{0\})\cap B(0,\varepsilon)$ tepat berkomponen \omterm{def:b3:topology:connected}{terhubung}
$2$ (yakni tudung tertusuk atas dan bawahnya, yang masing-masingnya
\omterm{def:b3:topology:pathconnected}{terhubung lintasan}: hubungkanlah lewat lingkaran dan sinarnya). Seandainya $C$
merupakan submanifold-$2$ di $0$, maka sebuah pelurusan akan memberikan
\omterm{def:b3:topology:continuity}{homeomorfisma} dari $C\cap\Omega$ ke sebuah kepingan terbuka sebuah
bidang yang mengirim $0$ ke sebuah titik $p$; padahal \omterm{def:b3:topology:topology}{persekitaran} bidang
tertusuk yang kecil di sekitar $p$ berkomponen \emph{satu}, dan
\omterm{def:b3:topology:continuity}{homeomorfisma} mengawetkan banyaknya komponen \omterm{def:b3:topology:topology}{persekitaran}
tertusuknya: jadi bertentangan.
\end{solution}

\begin{solution}{exo:b3:submanifolds:2}
(a) $T_aS^{n-1} = \ker\bigl(2a^{\mathsf T}\bigr) = a^\perp$.
(b) Di $p = ((R+r)\cos\theta, (R+r)\sin\theta, 0)$: $\nabla G
= (2r\cos\theta, 2r\sin\theta, 0)$, sehingga bidang singgungnya adalah
$\operatorname{Vect}\bigl((-\sin\theta, \cos\theta, 0),\ (0,
0, 1)\bigr)$: yakni bidang tegak yang menyinggung khatulistiwa
luarnya.
(c) Heliksnya merupakan kurva terbenam dengan $\varphi'(t_0) =
(-\sin t_0, \cos t_0, 1) \neq 0$: sehingga garis singgungnya di
$\varphi(t_0)$ adalah $\varphi(t_0) + \R\,\varphi'(t_0)$, seperti yang
ditetapkan \cref{prop:b3:submanifolds:tangent}(2).
\end{solution}

\begin{solution}{exo:b3:submanifolds:3}
(a) $Df(x,y) = \bigl(\begin{smallmatrix}2x & -2y\\ 2y &
2x\end{smallmatrix}\bigr)$ dengan $\det = 4(x^2 + y^2)$: sehingga
teoremanya berlaku di setiap titik kecuali titik asalnya.
(b) Karena $f(-z) = f(z)$: ia tak pernah injektif pada himpunan yang setangkup terhadap
$0$; sedangkan pada $\R^2\setminus\{0\}$ ia difeomorfisma lokal
di mana-mana namun $2$-ke-$1$ secara global. Lalu di dekat $(1, 0) = f(\pm(1,
0))$, kedua balikannya adalah kedua cabang akar kuadratnya: yakni dalam
notasi kompleksnya $w \mapsto \pm\sqrt w$ (dengan cabang pokoknya), yakni
\[
(u, v) \longmapsto \pm\Bigl(\sqrt{\tfrac{u +
\sqrt{u^2+v^2}}{2}},\ \frac{v}{2\sqrt{(u +
\sqrt{u^2+v^2})/2}}\Bigr).
\]
(c) Jacobinya $r > 0$: jadi difeomorfisma lokal pada
$\intoo0\infty\times\R$, tetapi $\theta \mapsto \theta + 2\pi$
memberikan titik yang sama: jadi terbalikkan secara lokal (sebab sudutnya tertentukan sampai
$2\pi$ pada sebuah setengah bidang), tetapi tak pernah secara global.
\end{solution}

\begin{solution}{exo:b3:submanifolds:4}
(a) Gradien $\nabla F = 3(x^2 - y,\ y^2 - x)$ lenyap jika dan hanya jika $y = x^2$
dan $x = y^2$, yakni $x^4 = x$: yaitu di $(0,0)$ dan $(1,1)$; sedangkan hanya
$(0,0)$ yang terletak pada $\mathcal C$ (sebab $F(1,1) = -1$). Jadi pada $\mathcal
C\setminus\{0\}$, $F$ merupakan submersi: yakni submanifold-$1$, yang secara
lokal berupa grafik $y = \psi(x)$ di tempat $F_y = 3(y^2 - x) \neq 0$
dan $x = \chi(y)$ di tempat $F_x = 3(x^2 - y) \neq 0$ (dan setidaknya
satu berlaku di luar titik asalnya).

(b) Di $(\frac32, \frac32)$: $\nabla F = 3(\frac94 - \frac32)
(1, 1) = \frac94(1,1)$: sehingga garis singgungnya $x + y = 3$.

(c) Parametrisasi rasionalnya $x = \frac{3t}{1 + t^3}$ dan
$y = \frac{3t^2}{1+t^3}$ melalui $0$ di $t = 0$ dengan
kecepatan $(3, 0)$; sedangkan menukar $x \leftrightarrow y$ (sebab
kurvanya setangkup, atau parametrikan ulang lewat $1/t$) memberikan kurva
$\mathcal C^1$ kedua yang melalui $0$ dengan kecepatan $(0, 3)$. Padahal dua
arah singgung yang bebas mustahil bagi sebuah
submanifold-$1$ (sebab \omterm{def:b3:submanifolds:tangent}{ruang singgungnya} sebuah garis,
\cref{prop:b3:submanifolds:tangent}): jadi $\mathcal C$ bukan
submanifold di titik asalnya --- yakni persilangan diri yang transversal.
\end{solution}

\begin{solution}{exo:b3:submanifolds:5}
(a) $F(M + H) = M^{\mathsf T}M + M^{\mathsf T}H + H^{\mathsf
T}M + H^{\mathsf T}H$: sehingga $DF(M)(H) = M^{\mathsf T}H +
H^{\mathsf T}M$. Lalu untuk $M \in O_n$ dan $S$ yang setangkup, $H =
\frac12MS$ memberikan $DF(M)(H) = \frac12(S + S^{\mathsf T}) = S$:
jadi surjektif ke $S_n$.

(b) $O_n = F^{-1}(I)$ dengan $F$ sebuah submersi (ke $S_n$ yang
berdimensi $\frac{n(n+1)}2$) di setiap titiknya:
jadi submanifold berdimensi $n^2 - \frac{n(n+1)}2 =
\frac{n(n-1)}2$. Dan \omterm{def:b3:topology:compact}{kompak}: sebab tertutup ($F$ \omterm{def:b3:topology:continuity}{kontinu}) dan terbatas
(karena kolomnya vektor satuan). Singgungnya di $I$: $\ker DF(I) =
\{H : H + H^{\mathsf T} = 0\}$.

(c) $\bigl(\eu^{tH}\bigr)^{\mathsf T}\eu^{tH} =
\eu^{tH^{\mathsf T}}\eu^{tH} = \eu^{-tH}\eu^{tH} = I$
(sebab transposnya lewat deretnya; dan \omterm{thm:b3:ode:matrixexp}{eksponensial matriks} yang
komutatif saling mengalikan,
\cref{thm:b3:ode:matrixexp}).
\end{solution}

\begin{solution}{exo:b3:submanifolds:6}
(a) $\det(M + H) = \det M\,\det(I + M^{-1}H) = \det M\bigl(1
+ \operatorname{tr}(M^{-1}H) + O(\norm H^2)\bigr)$ untuk $M$ yang
terbalikkan (lewat uraian $\det$ di dekat $I$: yakni suku linear
$\prod(1 + \lambda_i)$); lalu dengan $\det M\cdot M^{-1} =
\operatorname{com}(M)^{\mathsf T}$:
$D\det(M)(H) = \operatorname{tr}\bigl(\operatorname{com}
(M)^{\mathsf T}H\bigr)$, dan rumusnya diperluas ke setiap $M$
lewat \omterm{ex:b3:lebesgue:gamma}{kepadatan} dan \omterm{def:b3:topology:continuity}{kekontinuannya}. Sedangkan pada $SL_n$, $\det M = 1$:
sehingga $D\det(M) \ne 0$ (sebab nilainya pada $H = M$ adalah
$\operatorname{tr}(I)\cdot\dots = n\det M = n$).

(b) $SL_n = \det^{-1}(1)$ dengan $\det$ sebagai submersi di sana
(sebab nilainya di $\R$): jadi berdimensi $n^2 - 1$; dan $T_ISL_n = \ker
D\det(I) = \{H : \operatorname{tr}H = 0\}$.

(c) $GL_n$ merupakan himpunan bagian \emph{terbuka} $M_n(\R)$
(\cref{exo:b3:topology:8}): jadi submanifold berdimensi penuh
$n^2$ (dengan pelurusannya: yakni peta identitasnya).
\end{solution}

\begin{solution}{exo:b3:submanifolds:7}
(a) $(y, x) = \lambda(2x, 2y)$ dan $x^2 + y^2 = 1$: sehingga $y =
2\lambda x$ dan $x = 2\lambda y$ memberikan $x^2 = y^2 = \frac12$.
Nilai $xy$-nya: $\pm\frac12$: jadi maksimumnya $\frac12$ di
$\pm\frac1{\sqrt2}(1,1)$, dan minimumnya $-\frac12$ di
$\pm\frac1{\sqrt2}(1,-1)$ (sebab himpunan kendalanya \omterm{def:b3:topology:compact}{kompak}:
sehingga ekstremumnya ada).

(b) Pada \omterm{def:b3:topology:interior}{interior} simpleksnya (dengan $p_i > 0$), Lagrange bagi
$H(p) = -\sum p_i\ln p_i$ dengan kendala $\sum p_i = 1$:
$-\ln p_i - 1 = \lambda$ untuk setiap $i$: sehingga semua $p_i$ sama, $p_i =
\frac1n$, dengan $H = \ln n$. Lalu maksimum atas simpleks \omterm{def:b3:topology:compact}{kompaknya}
tercapai; dan seandainya ia tercapai pada batasnya
(dengan suatu $p_i = 0$), maka distribusinya hidup pada $\leq n - 1$
titik sehingga lewat induksi $H \leq \ln(n-1) < \ln n$: jadi
titik kritis \omterm{def:b3:topology:interior}{interiornya} adalah maksimum globalnya --- yakni
ketidaktahuan seragam yang memaksimumkan entropinya.

(c) Minimumkan $(x-1)^2 + y^2$ pada elips \omterm{def:b3:topology:compact}{kompaknya}:
$(2(x{-}1), 2y) = \lambda(\frac x2, 2y)$. Jika $y \neq 0$: maka
$\lambda = 1$, lalu $2(x - 1) = \frac x2$ memberikan $x =
\frac43$ dan $y^2 = 1 - \frac49 = \frac59$: sehingga jarak$^2$ $=
\frac19 + \frac59 = \frac23$. Sedangkan jika $y = 0$: $x = \pm2$,
dengan jarak $1$ dan $3$. Jadi titik terdekatnya: $\bigl(\frac43,
\pm\frac{\sqrt5}3\bigr)$, pada jarak $\sqrt{2/3} < 1$. Sedangkan
persamaan pengalinya mengatakan bahwa ruas dari $(1,0)$ ke
titik terdekatnya sejajar dengan $\nabla$(elipsnya): sehingga ia memotong
elipsnya secara ortogonal, seperti yang dituntut geometrinya.
\end{solution}

\begin{solution}{exo:b3:submanifolds:8}
Lewat induksi pada $n$; dan $n = 1$ trivial. Fungsi Rayleigh
$f(x) = \langle Ax, x\rangle$ mencapai maksimumnya $\lambda_1$
pada $S^{n-1}$ yang \omterm{def:b3:topology:compact}{kompak} di suatu $v_1$; lalu Lagrange
(\cref{thm:b3:submanifolds:lagrange}, dengan bolanya sebagai himpunan aras)
memberikan $2Av_1 = 2\lambda v_1$, dan $\lambda = \langle Av_1,
v_1\rangle = \lambda_1$. Sedangkan hiperbidang $v_1^\perp$ bersifat
invarian-$A$ (menurut kesetangkupannya: $\langle Av, v_1\rangle = \langle
v, Av_1\rangle = 0$); pembatasannya setangkup, dan
induksinya melahirkan basis eigen ortonormal $v_2, \dots, v_n$
milik $v_1^\perp$ bernilai eigen $\lambda_2 \geq \dots \geq
\lambda_n$, yang masing-masingnya maksimum $f$ pada bola
\omterm{thm:b3:hilbert:decomposition}{komplemen ortogonal} sisanya. Lalu Courant--Fischer menyusul persis
seperti pada \cref{exo:b3:spectral:8}: uraikan $x = \sum c_iv_i$;
lalu pada sebuah ruang uji berdimensi $k$ iriskanlah dengan
$\operatorname{Vect}(v_k, \dots, v_n)$ (lewat pencacahan dimensi di
$\R^n$) untuk memperoleh $\min \leq \lambda_k$, sedangkan
$\operatorname{Vect}(v_1, \dots, v_k)$ mencapai $\min =
\lambda_k$.
\end{solution}

\begin{solution}{exo:b3:submanifolds:9}
Menskalakan setiap kolomnya menjadi bernorma satuan membagi $\abs{\det}$ dengan
$\prod\norm{c_i}$: sehingga cukup dibuktikan $\abs{\det M} \leq 1$
bila semua kolomnya satuan, dengan kesamaannya jika dan hanya jika $M \in O_n$. Lalu
fungsi $\det$ \omterm{def:b3:topology:continuity}{kontinu} pada
$(S^{n-1})^n$ yang \omterm{def:b3:topology:compact}{kompak}: sehingga ia mencapai maksimum $m \geq \det I = 1 > 0$
di suatu $M$. Lalu dengan menetapkan semua kolomnya kecuali yang ke-$i$, $\det$ bersifat
linear terhadap $c_i$ dengan gradien berupa kolom ke-$i$
$\operatorname{com}(M)$; sehingga Lagrange pada bola ke-$i$-nya:
$\operatorname{com}(M)_{\cdot i} = \lambda_i c_i$. Sedangkan identitas
$M^{\mathsf T}\operatorname{com}(M) = \det(M)\,I$
berbunyi $\langle c_j, \operatorname{com}(M)_{\cdot i}\rangle =
\det(M)\,\delta_{ij}$, yakni $\lambda_i\langle c_j,
c_i\rangle = \det(M)\delta_{ij}$; lalu dengan mengambil $j = i$: $\lambda_i
= \det M = m \neq 0$, dan lalu $j \neq i$ memberikan $\langle c_i,
c_j\rangle = 0$: sehingga kolomnya ortonormal, $M \in O_n$, dan
$m = \abs{\det M} = 1$. Karena itu $\abs{\det} \leq
\prod\norm{c_i}$ selalu, dengan kesamaannya tepat bagi kolom yang
ortogonal (skalakan kembali): sehingga volume sebuah paralelepipedum terbesar,
bagi panjang rusuk yang diberikan, saat rusuknya saling tegak lurus.
\end{solution}

\begin{solution}{exo:b3:submanifolds:10}
Grafik $y = \pm\sqrt{1 - x^2}$ berhasil di dekat setiap titik yang $y \neq
0$; sedangkan $x = \pm\sqrt{1 - y^2}$ di dekat setiap titik yang $x \neq 0$;
dan di $(1, 0)$: grafik $x = \sqrt{1 - y^2}$ atas $y \in
\intoo{-1}1$, yang merupakan
\cref{thm:b3:submanifolds:characterizations}(3) dengan
koordinatnya ditukar. Sedangkan suatu permutasi selalu berhasil karena
garis singgungnya, yang berdimensi satu, tak dapat sekaligus
tegak dan mendatar --- tetapi ia dapat menjadi salah satunya, sehingga
tak ada pilihan koordinat ``terikat'' yang tetap yang
melayani setiap titiknya.
\end{solution}

\begin{solution}{exo:b3:submanifolds:11}
(a) Pemetaan pendefinisinya $F(M) = M^{\mathsf T}M - I$ bersifat
\omterm{def:b3:topology:continuity}{kontinu}: sehingga $O_n = F^{-1}(0)$ tertutup; sedangkan kolom sebuah matriks
ortogonal berupa vektor satuan, sehingga norma Euclidnya
(yakni \omterm{thm:b3:galois:finitefields}{Frobenius}) tepat $\sqrt n$: jadi terbatas. Sehingga \omterm{def:b3:topology:compact}{kompak} menurut
Heine--Borel di $M_n(\R) \cong \R^{n^2}$.

(b) Himpunan $O_n$ adalah himpunan aras $F = 0$ yang dipelajari pada babnya:
$DF(A)H = A^{\mathsf T}H + H^{\mathsf T}A$, yang surjektif ke
matriks setangkupnya di setiap $A \in O_n$ (sebab diberikan $S$ yang setangkup,
ambillah $H = \frac12AS$), sehingga $O_n$ merupakan submanifold
berdimensi $n^2 - \frac{n(n+1)}2 = \frac{n(n-1)}2$ dengan
\[
T_AO_n = \ker DF(A) = \{H : A^{\mathsf T}H \text{
antisetangkup}\} = \{AK : K^{\mathsf T} = -K\} ;
\]
dan di $A = I$ semuanya berupa matriks antisetangkup.

(c) $\exp(tK)^{\mathsf T}\exp(tK) = \exp(tK^{\mathsf T})
\exp(tK) = \exp(-tK)\exp(tK) = I$ (sebab kedua matriks $\pm tK$
komutatif, sehingga hasil kali eksponensialnya adalah eksponensial
jumlahnya): jadi $\exp(tK) \in O_n$, dan $\gamma(t) =
A\exp(tK)$ merupakan kurva di $O_n$ dengan $\gamma(0) = A$ dan
$\gamma'(0) = AK$. Lalu saat $K$ menjelajah matriks antisetangkupnya,
$AK$ menyapu $T_AO_n$: sehingga eksponensialnya mewujudkan seluruh
\omterm{def:b3:submanifolds:tangent}{ruang singgungnya} lewat kurva yang eksplisit --- yakni jalan pintas grup Lie
yang dimanfaatkan \cref{pb:b3:submanifolds:1} bagi $SO(3)$.
\end{solution}

\begin{solution}{exo:b3:submanifolds:12}
Untuk keberadaannya: iriskan $M$ dengan sebuah bola tertutup besar di sekitar
$p$ untuk memperoleh sebuah \omterm{def:b3:topology:compact}{kompak} tak kosong; lalu jaraknya yang \omterm{def:b3:topology:continuity}{kontinu}
mencapai minimumnya di sana, sedangkan titik di luar bolanya lebih
jauh. Untuk syarat orde pertamanya: bagi sebuah kurva $\gamma$ di $M$
dengan $\gamma(0) = x_0$, fungsi $h(t) = \norm{\gamma(t)
- p}^2$ dapat diturunkan dengan minimum di $0$:
\[
0 = h'(0) = 2\,\langle\gamma'(0),\ x_0 - p\rangle,
\]
dan $\gamma'(0)$ menyapu $T_{x_0}M$: sehingga $p - x_0 \perp
T_{x_0}M$. Untuk bolanya $S(c, r)$: \omterm{def:b3:submanifolds:tangent}{ruang singgungnya} di $x_0$ adalah
$(x_0 - c)^\perp$, sehingga $p - x_0 \parallel x_0 - c$: jadi $x_0$
terletak pada garis yang melalui $c$ dan $p$, berjarak $r$ dari
$c$ --- yakni titik sinarnya, seperti yang didesak geometrinya. Untuk parabolanya: di
$x_0 = (x, x^2)$ singgungnya direntang $(1, 2x)$;
sehingga keortogonalannya terhadap $p - x_0 = (2 - x, -x^2)$ berbunyi
\[
(2 - x) - 2x^3 = 0,
\qquad\text{yakni}\qquad 2x^3 + x - 2 = 0,
\]
dengan akar real tunggalnya (sebab $x \mapsto 2x^3 + x$ naik
sejati) $x \approx 0.835$; lalu $x_0 \approx (0.835,
0.698)$ dan $d(p, M) = \sqrt{(2 - 0.835)^2 + 0.698^2}
\approx 1.358$.
\end{solution}

\begin{solution}{pb:b3:submanifolds:1}
\textbf{1.} Petakan $q = t + x\mathrm i + y\mathrm j + z\mathrm k
\mapsto \bigl(\begin{smallmatrix}\alpha & \beta\\ -\bar\beta
& \bar\alpha\end{smallmatrix}\bigr)$ dengan $\alpha = t + \iu
x$ dan $\beta = y + \iu z$: lalu diperiksa bahwa $1, \mathrm i,
\mathrm j, \mathrm k$ berpindah ke $I$, $\bigl(\begin{smallmatrix}
\iu & 0\\ 0 & -\iu\end{smallmatrix}\bigr)$,
$\bigl(\begin{smallmatrix}0 & 1\\ -1 &
0\end{smallmatrix}\bigr)$, dan $\bigl(\begin{smallmatrix}0 &
\iu\\ \iu & 0\end{smallmatrix}\bigr)$, yang hasil kalinya
menghasilkan kembali tabel kuaternionnya: sehingga pemetaannya merupakan morfisma
aljabar yang injektif, jadi $\mathbb H$ mewarisi keasosiatifannya;
lalu $N(q) = \abs\alpha^2 + \abs\beta^2 = \det$ bersifat multiplikatif,
dan pengonjugatannya bersesuaian dengan transpos adjugatnya,
yang memberikan $\overline{pq} = \bar q\bar p$. Untuk pusatnya: komutatif
dengan $\mathrm i$ memaksa $y = z = 0$, dan dengan $\mathrm j$ memaksa
$x = 0$: jadi $\R$.

\textbf{2.} Karena $q\bar q = N(q)$: maka untuk $q \neq 0$, $q^{-1} = \bar
q/N(q)$: jadi sebuah aljabar pembagian. Sedangkan pada $S^3$: $N(pq) = 1$ dan
$N(q^{-1}) = 1$: jadi sebuah grup; dan $S^3 \subseteq \R^4$ adalah
bola satuannya: jadi submanifold-$3$ yang \omterm{def:b3:topology:compact}{kompak}.

\textbf{3.} Sebuah $v$ bersifat murni jika dan hanya jika $\bar v = -v$; lalu
$\overline{qv\bar q} = q\bar v\bar q = -qv\bar q$: sehingga $\rho_q$
mengawetkan $P$. Kelinearannya jelas; sedangkan $N(qv\bar q) =
N(q)N(v)N(q) = N(v)$: jadi sebuah isometri $(P, N) \cong (\R^3,
\norm\cdot^2)$: sehingga $\rho_q \in O(3)$. Dan $\rho_{pq}(v) =
pqv\overline{pq} = p(qv\bar q)\bar p =
\rho_p(\rho_q(v))$: jadi sebuah morfisma.

\textbf{4.} Berlaku $\rho_q = \mathrm{id}$ jika dan hanya jika $qv = vq$ untuk setiap $v$
yang murni, jika dan hanya jika $q$ komutatif dengan $\mathrm i, \mathrm j, \mathrm k$,
jika dan hanya jika $q$ bersifat pusat (pertanyaan 1): jadi $q \in \R\cap S^3 =
\{\pm1\}$.

\textbf{5.} Tulis $q = t + p$ (dengan $t \in \R$ dan $p$ murni): maka $1 =
N(q) = t^2 + N(p)$, sehingga $t = \cos\frac\theta2$ dan $p =
\sin\frac\theta2\,u$ dengan $N(u) = 1$ untuk suatu $\theta$
(dan jika $p = 0$, maka $q = \pm1$ bekerja secara trivial). Lalu karena $u^2 = -N(u)
= -1$, $q$ dan $u$ komutatif, sehingga $\rho_q(u) = qu\bar q =
uq\bar q = u$: yakni sumbunya. Sedangkan untuk satuan murni $w \perp u$:
aturan hasil kalinya $vw = -\langle v, w\rangle + v\wedge w$
(uraikanlah dalam koordinatnya dari tabelnya) memberikan $uw = u\wedge w
= -wu$. Maka
\[
\rho_q(w) = \bigl(\cos\tfrac\theta2 +
\sin\tfrac\theta2u\bigr)\,w\,\bigl(\cos\tfrac\theta2 -
\sin\tfrac\theta2u\bigr)
= \cos\theta\,w + \sin\theta\,(u\wedge w),
\]
dengan memakai $uwu = -u^2w = w$ dan rumus sudut gandanya: yakni
rotasi bersudut $\theta$ pada bidang berarah $(w,
u\wedge w)$.

\textbf{6.} Setiap $\rho_q$ merupakan rotasi terhadap $u$ sebesar
$\theta$: sebab pada basis ortonormal $(u, w, u\wedge w)$
matriksnya berdeterminan $+1$: jadi $\operatorname{im}\rho \subseteq
SO(3)$. Untuk kesurjektifannya: sebuah matriks $R \in SO(3)$ bernilai eigen $1$, sebab
\[
\det(R - I) = \det R\,\det(I - R^{\mathsf T}) = \det(I - R)
= (-1)^3\det(R - I),
\]
sehingga $\det(R - I) = 0$. Lalu ambillah sebuah vektor eigen satuan $u$; maka $R$
mengawetkan $u^\perp$ dan terbatas di sana menjadi sebuah rotasi bersudut
suatu $\theta$ (sebab ortogonal bidang berdeterminan $1$):
sehingga $R = \rho_q$ untuk $q = \cos\frac\theta2 +
\sin\frac\theta2\,u$. Lalu dengan pertanyaan 4 dan teorema
isomorfisma pertamanya (\cref{thm:b3:groups:firstiso}): $SO(3)
\cong S^3/\{\pm1\}$.

\textbf{7.} Himpunan $O_3$ merupakan submanifold \omterm{def:b3:topology:compact}{kompak} berdimensi $3$
(\cref{exo:b3:submanifolds:5}); sedangkan $\det$ \omterm{def:b3:topology:continuity}{kontinu} padanya
dengan nilai di $\{\pm1\}$, sehingga $SO(3) = O_3\cap\{\det = 1\}$
bersifat terbuka sekaligus tertutup di $O_3$: jadi sebuah gabungan \omterm{def:b3:topology:components}{komponen terhubungnya},
sehingga ia sendiri submanifold-$3$ yang \omterm{def:b3:topology:compact}{kompak}, dengan \omterm{def:b3:submanifolds:tangent}{ruang singgung}
yang sama di $I$: yakni matriks antisetangkupnya.

\textbf{8.} Berlaku $A_u^2v = u\wedge(u\wedge v) = \langle u,
v\rangle u - v$ (untuk $u$ yang satuan), sehingga $A_u^3v = u\wedge(\langle
u,v\rangle u - v) = -u\wedge v$: jadi $A_u^3 = -A_u$. Lalu dengan memisahkan
deret eksponensialnya menurut sisa pangkatnya modulo relasi
$A^3 = -A$:
\[
\eu^{\theta A_u} = I + \Bigl(\theta - \frac{\theta^3}{3!} +
\cdots\Bigr)A_u + \Bigl(\frac{\theta^2}{2!} -
\frac{\theta^4}{4!} + \cdots\Bigr)A_u^2
= I + \sin\theta\,A_u + (1 - \cos\theta)\,A_u^2 .
\]
Pada $u$: $A_uu = 0$: jadi tetap. Sedangkan pada $w \perp u$: $\eu^{\theta
A_u}w = w + \sin\theta\,u\wedge w + (1 - \cos\theta)(-w) =
\cos\theta\,w + \sin\theta\,u\wedge w$: yakni rotasi bersumbu
$u$ dan bersudut $\theta$ --- yaitu Rodrigues. Dan setiap rotasi berbentuk
demikian (pertanyaan 6): sehingga $\exp$ bersifat pada $SO(3)$ dari
matriks antisetangkupnya.

\textbf{9.} Dengan $q(t) = \cos\frac t2 + \sin\frac t2\,u$:
pertanyaan 5 menunjukkan bahwa $\rho_{q(t)}$ adalah rotasi bersumbu $u$
dan bersudut $t$, yakni $\rho_{q(t)} = \eu^{tA_u}$, yang
turunannya di $t = 0$ adalah $A_u$. Jadi kuaternionnya berjalan pada separuh
sudutnya --- yakni jejak analitis selimut gandanya.

\textbf{10.} Pemetaan $\rho_{q(t)}$ adalah rotasi terhadap $\mathrm k$
sebesar sudut $t$: sehingga di $t = 2\pi$ ia kembali ke identitasnya ---
yakni sebuah gelung di $SO(3)$. Sedangkan angkatannya memenuhi $q(2\pi) = \cos\pi =
-1 = -q(0)$: jadi \omterm{def:b3:topology:pathconnected}{lintasan} terangkatnya \emph{tak} tertutup; dan hanya di
$t = 4\pi$ $q$ kembali ke $1$. Tafsirannya: bahwa gelung
rotasi penuhnya tak terkontraksikan di $SO(3)$ --- sebab
angkatannya berakhir di lembar lain selimutnya --- sedangkan
gelung gandanya terkontraksikan; sehingga sebuah benda yang terikat ke sekitarnya oleh
tali (yakni kiat sabuknya) kembali ke keadaan tak terpilin setelah
$4\pi$ tetapi tidak setelah $2\pi$. Jadi grup rotasinya mengingat
paritas putaran penuhnya; dan $S^3$, yang \omterm{def:b3:conformal:simplyconnected}{terhubung sederhana}, adalah
tempat ingatan itu tinggal.

\textbf{11.} Untuk seperempat putarannya: $q_{\mathrm i} =
\cos\frac\pi4 + \sin\frac\pi4\,\mathrm i$ dan $q_{\mathrm j} =
\cos\frac\pi4 + \sin\frac\pi4\,\mathrm j$. Lalu hasil kalinya
(dengan putaran-$\mathrm j$ diterapkan lebih dahulu):
\[
q_{\mathrm i}q_{\mathrm j} = \tfrac12(1 + \mathrm i)(1 +
\mathrm j) = \tfrac12\bigl(1 + \mathrm i + \mathrm j +
\mathrm k\bigr),
\]
yang bernorma $1$, dengan $\cos\frac\theta2 = \frac12$: sehingga $\theta =
\frac{2\pi}3$, dan sumbunya $u = \frac{\mathrm i + \mathrm j +
\mathrm k}{\sqrt3}$ (yakni bagian murninya yang dinormalkan). Jadi dua
seperempat putaran berurutan terhadap sumbu yang ortogonal menghasilkan
rotasi $120^\circ$ terhadap diagonal utama kubusnya --- yakni pembukuan
senilai empat perkalian real, dengan keortogonalannya
terawetkan secara persis: itulah mengapa perangkat lunak penerbangan dan mesin grafis
mengomposisikan rotasinya lewat kuaternion.

\textbf{12.} Dari tabelnya, $\mathrm{ji} = -\mathrm k$ dan
$\mathrm{ki} = \mathrm j$, sehingga
\[
q\,\mathrm i = a\mathrm i - b + c(\mathrm{ji}) +
d(\mathrm{ki}) = -b + a\mathrm i + d\mathrm j - c\mathrm k .
\]
Lalu mengalikannya dengan $\bar q = a - b\mathrm i - c\mathrm j -
d\mathrm k$ memakai aturan skalar dan vektornya $(t_1 + p_1)(t_2 +
p_2) = t_1t_2 - \langle p_1, p_2\rangle + t_1p_2 + t_2p_1 +
p_1\wedge p_2$, dengan $p_1 = (a, d, -c)$ dan $p_2 = (-b, -c,
-d)$: bagian skalarnya adalah $-ab - (-ab) = 0$ (jadi murni, seperti yang
seharusnya), dan bagian vektornya adalah
\[
b(b, c, d) + a(a, d, -c) + (-c^2 - d^2,\ bc + ad,\ bd - ac)
= \bigl(a^2 + b^2 - c^2 - d^2,\ 2(bc + ad),\ 2(bd -
ac)\bigr):
\]
yakni kolom pertama $R_q$. Sedangkan pemetaan siklik $\sigma(\mathrm i)
= \mathrm j$, $\sigma(\mathrm j) = \mathrm k$, dan
$\sigma(\mathrm k) = \mathrm i$ mengawetkan relasi
$\mathrm i^2 = \mathrm j^2 = \mathrm k^2 = \mathrm{ijk} =
-1$ (sebab kata $\mathrm{ijk}$ invarian secara siklik sampai
relasi $\mathrm{ijk} = \mathrm{jki}$, yang berlaku di
ring mana pun: dengan mengonjugatkan $\mathrm{ijk} = -1$ oleh
$\mathrm i$ yang terbalikkan), sehingga $\sigma$ diperluas menjadi
automorfisma aljabar-$\R$, dan $\sigma(\rho_q(v)) =
\rho_{\sigma(q)}(\sigma(v))$. Lalu dengan menguraikannya, peta
$\mathrm j$ adalah rumus kolom pertamanya setelah
substitusi $(b, c, d) \to (c, d, b)$ dengan basisnya
dilabeli ulang $\mathrm i \to \mathrm j \to \mathrm k \to \mathrm
i$, yang tepat merupakan kolom kedua yang ditampilkan; dan satu
putaran lagi memberikan yang ketiga.

\textbf{13.} Dengan menjumlahkan diagonalnya, $\operatorname{tr}R_q =
3a^2 - (b^2 + c^2 + d^2) = 4a^2 - 1$ (menurut norma satuannya), dan dengan
$a = \cos\frac\theta2$: $4\cos^2\frac\theta2 - 1 = 1 +
2\cos\theta$. Untuk bagian antisetangkupnya: ketiga entri bebas
$R_q - R_q^{\mathsf T}$ adalah $4ab, 4ac, 4ad$ (pada
posisi $(3,2), (1,3), (2,1)$), sehingga $\frac12(R_q -
R_q^{\mathsf T}) = A_m$ dengan $m = 2a\,(b, c, d) =
2\cos\frac\theta2\sin\frac\theta2\,u = \sin\theta\,u$.
Algoritmanya: $\theta = \arccos\frac{\operatorname{tr}R -
1}{2} \in \intcc0\pi$; lalu jika $0 < \theta < \pi$, bacalah $u$ dari
$\frac{R - R^{\mathsf T}}{2\sin\theta}$ lalu tetapkan $q =
\pm(\cos\frac\theta2 + \sin\frac\theta2 u)$; sedangkan jika $\theta =
0$, maka $q = \pm1$. Untuk $\theta = \pi$: $a = 0$, dan Rodrigues
(pertanyaan 8) memberikan $R = I + 2A_u^2 = I + 2(uu^{\mathsf T} -
I) = 2uu^{\mathsf T} - I$, yakni $R + I = 2uu^{\mathsf T}$;
sehingga sembarang kolom tak nol $R + I$, yang dinormalkan, adalah $\pm u$, dan
$q = \pm u$.

\textbf{14.} Dengan $a = b = c = d = \frac12$: semua entri diagonalnya
lenyap, $2(bc - ad) = 0$, $2(bd + ac) = 1$, $2(bc +
ad) = 1$, $2(cd - ab) = 0$, $2(bd - ac) = 0$, dan $2(cd + ab) =
1$:
\[
R_q = \begin{pmatrix} 0 & 0 & 1\\ 1 & 0 & 0\\ 0 & 1 & 0
\end{pmatrix},
\]
yakni permutasi siklik $e_1 \to e_2 \to e_3 \to e_1$. Sedangkan
tracenya $0 = 1 + 2\cos\theta$, sehingga $\theta = \frac{2\pi}3$,
dan ia menetapkan $(1,1,1)$: yakni rotasi sebesar $120^\circ$ terhadap
diagonal utamanya --- tepat jawaban pertanyaan 11, yang kini
terlihat sebagai matriks yang menyiklikkan sumbu koordinatnya.

\textbf{15.} Pada basisnya diperiksa bahwa $\Phi(\bar q) =
\Phi(q)^*$ (sebab matriks $\bar q$ mempunyai $\alpha' =
\bar\alpha$ dan $\beta' = -\beta$, yang merupakan transpos
konjugat $\bigl(\begin{smallmatrix}\alpha & \beta\\
-\bar\beta & \bar\alpha\end{smallmatrix}\bigr)$). Karena itu
$\Phi(q)^*\Phi(q) = \Phi(\bar qq) = N(q)I$ dan $\det\Phi(q)
= \abs\alpha^2 + \abs\beta^2 = N(q)$: sehingga untuk $q \in S^3$,
$\Phi(q) \in SU(2)$, dan $\Phi$ merupakan morfisma yang injektif
(pertanyaan 1). Untuk kesurjektifannya: misalkan $U =
\bigl(\begin{smallmatrix}\alpha & \beta\\ \gamma &
\delta\end{smallmatrix}\bigr)$ dengan $\det U = 1$; maka
$U^{-1} = \bigl(\begin{smallmatrix}\delta & -\beta\\
-\gamma & \alpha\end{smallmatrix}\bigr)$, dan $U^{-1} = U^*
= \bigl(\begin{smallmatrix}\bar\alpha & \bar\gamma\\
\bar\beta & \bar\delta\end{smallmatrix}\bigr)$ memaksa
$\delta = \bar\alpha$ dan $\gamma = -\bar\beta$, lalu $1 =
\det U = \abs\alpha^2 + \abs\beta^2$: sehingga $U = \Phi(q)$ untuk
kuaternion satuan $q$ berkoordinat $\alpha = a + \iu b$ dan
$\beta = c + \iu d$. Jadi $S^3 \cong SU(2)$.

\textbf{16.} Skalar realnya bersifat pusat dan $N(p) = 1$ memberikan
$\overline{pq\bar p} = p\bar q\bar p$, sehingga $pq\bar p +
\overline{pq\bar p} = p(q + \bar q)\bar p = q + \bar q$:
jadi bagian realnya invarian. Sebaliknya misalkan
$\operatorname{Re}q = \operatorname{Re}q' = a$; maka
bagian murninya bernorma sama $\sqrt{1 - a^2} = s$. Jika $s =
0$, maka $q = q' = \pm1$. Sedangkan jika $s > 0$, tulislah $q = a + su$ dan $q' =
a + su'$ dengan $u, u'$ murni satuan; lalu pertanyaan 6 menyediakan sebuah
rotasi yang membawa $u'$ ke $u$, yakni $p \in S^3$ dengan
$\rho_p(u') = u$, sehingga $pq'\bar p = a + s\rho_p(u') =
q$. Jadi kelas $S^3$-nya adalah $\{1\}$, $\{-1\}$,
dan untuk setiap $a \in \intoo{-1}1$ bola-$2$ $\{a + su :
u \text{ murni satuan}\}$ berjari-jari $s$. Lalu di bawah $\Phi$,
$\operatorname{tr}\Phi(q) = \alpha + \bar\alpha =
2\operatorname{Re}q$: sehingga kelas $SU(2)$-nya adalah himpunan
aras tracenya. Lalu dengan memproyeksikannya: jika $q' = pq\bar p$ maka
$\rho_{q'} = \rho_p\rho_q\rho_p^{-1}$; sebaliknya
$\rho_{q'} = \rho_p\rho_q\rho_p^{-1} = \rho_{pq\bar p}$
memaksa $q' = \pm pq\bar p$ (lewat kernelnya), sehingga
$\operatorname{Re}q' = \pm\operatorname{Re}q$, yakni
$\cos\frac{\theta'}2 = \abs{\operatorname{Re}q'} =
\abs{\operatorname{Re}q} = \cos\frac\theta2$ bagi sudutnya
di $\intcc0\pi$: jadi rotasi yang konjugat bersudut sama.
Sebaliknya, sudut yang sama mengizinkan wakil yang berbagian
real tak negatif sama, yang konjugat lewat cara di atas: jadi di
$SO(3)$, \omterm{ex:b3:groups:actions}{kelas konjugasi} sebuah rotasi tepat berupa
sudutnya.

\textbf{17.} Karena $N$ multiplikatif, $\abs\cdot$ merupakan norma
multiplikatif pada $\mathbb H \cong \R^4$ dan
$\abs{q^k} = \abs q^k$: sehingga deret $\sum q^k/k!$ konvergen
mutlak pada ruang berdimensi berhingga (jadi \omterm{def:b3:complete:complete}{lengkap})
itu, yang terdominasi $\sum\abs q^k/k! = \eu^{\abs q}$. Lalu untuk
$u$ yang murni satuan: $u^2 = -1$, sehingga $(\theta u)^{2m} =
(-1)^m\theta^{2m}$ dan $(\theta u)^{2m+1} =
(-1)^m\theta^{2m+1}u$; lalu dengan memisahkan deretnya,
\[
\exp(\theta u) = \sum_m\frac{(-1)^m\theta^{2m}}{(2m)!} +
u\sum_m\frac{(-1)^m\theta^{2m+1}}{(2m+1)!} = \cos\theta +
\sin\theta\,u .
\]
Sedangkan setiap $q \in S^3$ berupa $\cos\alpha + \sin\alpha\,u$ dengan
$\alpha \in \intcc0\pi$ (pertanyaan 5): jadi $q = \exp(\alpha u)$,
sehingga $\exp(P) = S^3$. Akhirnya $\exp(su) = \cos s + \sin s\,u
= q(2s)$ dalam notasi pertanyaan 9, dan $\rho_{q(t)} =
\eu^{tA_u}$ di sana: sehingga $\rho_{\exp(su)} = \eu^{2sA_u}$.

\textbf{18.} Dengan menguraikan $vw$ koordinat demi koordinat lewat tabelnya:
hasil kali $\mathrm i\cdot\mathrm i = -1$, \dots\ memberikan
skalar $-(v_1w_1 + v_2w_2 + v_3w_3)$, sedangkan hasil kali
campurannya ($\mathrm{ij} = \mathrm k$, $\mathrm{ji} =
-\mathrm k$, \dots) memberikan vektor $(v_2w_3 - v_3w_2,\
v_3w_1 - v_1w_3,\ v_1w_2 - v_2w_1)$: jadi $vw = -\langle v,
w\rangle + v\wedge w$. Lalu dengan mengurangkan hasil kali terbaliknya:
$vw - wv = 2\,v\wedge w$ (sebab bagian skalarnya meniadakan diri, sedangkan
hasil kali silangnya menjumlah). Untuk matriksnya, dengan $a\wedge(b\wedge
c) = b\langle a, c\rangle - c\langle a, b\rangle$:
\[
[A_v, A_w]x = v\wedge(w\wedge x) - w\wedge(v\wedge x)
= w\langle v, x\rangle - v\langle w, x\rangle
= (v\wedge w)\wedge x = A_{v\wedge w}x .
\]
Untuk turunannya: $\overline{\exp(tv)} = \exp(-tv)$ (sebab pengonjugatannya
\omterm{def:b3:topology:continuity}{kontinu} dan menegasikan kuaternion murni), sehingga
\[
\frac{\dd}{\dd t}\Bigr|_{t=0}\exp(tv)\,x\,\exp(-tv) = vx -
xv = 2\,v\wedge x = 2A_vx :
\]
sehingga diferensialnya adalah $v \mapsto 2A_v$, yakni bijeksi linear
dari $P$ ke matriks antisetangkupnya, dan $[2A_v, 2A_w]
= 4A_{v\wedge w} = 2A_{2v\wedge w} = 2A_{[v,w]}$ menunjukkan bahwa ia
membawa \omterm{def:b3:groups:derived}{komutator} kuaternionnya ke \omterm{def:b3:groups:derived}{komutator}
matriksnya.

\textbf{19.} Dengan menerapkan $\rho$: $\rho(\varepsilon(t)) =
\rho(s(R(t)))\,\rho(q(t))^{-1} = R(t)R(t)^{-1} =
\operatorname{id}$, sehingga $\varepsilon(t) \in \ker\rho =
\{\pm1\}$ (pertanyaan 4). Lalu sebagai hasil kali \omterm{def:b3:topology:continuity}{pemetaan kontinu}
$t \mapsto s(R(t))$ dan $t \mapsto q(t)^{-1} =
\bar{q(t)}$, $\varepsilon$ bersifat \omterm{def:b3:topology:continuity}{kontinu} pada selang \omterm{def:b3:topology:connected}{terhubung}
$\intcc0{2\pi}$ dengan nilai pada pasangan diskret
$\{\pm1\}$: sehingga ia konstan, katakanlah $\varepsilon(t) \equiv
\varepsilon$. Padahal $R(0) = R(2\pi) = I$, sehingga $s(R(0)) =
s(R(2\pi))$, sedangkan $s(R(0)) = \varepsilon\,q(0) =
\varepsilon$ dan $s(R(2\pi)) = \varepsilon\,q(2\pi) =
-\varepsilon$: jadi bertentangan. Sehingga tak ada seksi \omterm{def:b3:topology:continuity}{kontinu}
yang ada: jadi ketaktentuan tanda $\pm q$ bersifat global, bukan cacat
sebuah rumus tertentu.

\textbf{20.} \emph{Untuk kesurjektifannya:} setiap $R \in SO(3)$ berupa
$\eu^{\theta A_u}$ untuk suatu $u$ satuan dan $\theta \in
\intcc0{2\pi}$ (pertanyaan 6 dan 8); lalu jika $\theta > \pi$,
Rodrigues memberikan $\eu^{\theta A_u} = \eu^{(2\pi -
\theta)A_{-u}}$ (sebab keduanya sama dengan $I + \sin\theta A_u + (1 -
\cos\theta)A_u^2$, dan $A_{-u} = -A_u$ dengan $\sin(2\pi -
\theta) = -\sin\theta$ serta $\cos(2\pi - \theta) =
\cos\theta$), sehingga $R = E(v)$ dengan $\norm v \leq \pi$.
\emph{Untuk keinjektifannya di dalam:} jika $E(v) = E(v') \neq I$ dengan
$\norm v, \norm{v'} < \pi$, maka pertanyaan 13 memulihkan sudut
$\theta = \norm v = \norm{v'} \in \intoo0\pi$ yang sama dari
tracenya dan, karena $\sin\theta \neq 0$, sumbu yang sama
dari bagian antisetangkupnya: jadi $v = v'$; sedangkan $E(v) = I$
memaksa $\theta \in \{0\}$ pada bola terbukanya. \emph{Untuk batasnya:}
$E(\pi u) = I + 2A_u^2 = 2uu^{\mathsf T} - I$ bergantung pada
$u$ hanya lewat $uu^{\mathsf T}$, sehingga $E(\pi u) =
E(-\pi u)$; dan sebaliknya $2uu^{\mathsf T} - I =
2u'u'^{\mathsf T} - I$ yang diterapkan pada $u$ memberikan $u = \langle
u', u\rangle u'$, sehingga $u' = \pm u$. Sedangkan \omterm{def:b3:topology:interior}{interior} dan batasnya
tak pernah bertumbukan (sebab tracenya $> -1$ lawan $= -1$). Jadi $E$ menginduksi
bijeksi \omterm{def:b3:topology:continuity}{kontinu} dari bola-dengan-perekatan-antipodalnya
--- yang \omterm{def:b3:topology:compact}{kompak} --- ke $SO(3)$: yakni sebuah \omterm{def:b3:topology:continuity}{homeomorfisma}, sehingga $SO(3)
\cong \mathbb{RP}^3$. Dan sebuah diameter dari $\pi u$ ke $-\pi u$
berujung terekat: jadi ia sebuah gelung di $SO(3)$, dan
pemerian-$E$-nya cocok dengan keluarga rotasi pertanyaan 10
terhadap $u$ yang menyapu satu putaran penuh.

\textbf{21.} Karena $\rho$ merupakan morfisma dan $\rho_p^{-1} =
\rho_{p^{-1}} = \rho_{\bar p}$, maka $\rho_p\rho_q\rho_p^{-1}
= \rho_{pq\bar p}$; lalu menurut pertanyaan 16, mengonjugatkan sebuah rotasi
mengawetkan sudutnya dan memutar sumbunya oleh $\rho_p$.
Untuk involusinya: $\rho_q^2 = \rho_{q^2} = \operatorname{id}$
jika dan hanya jika $q^2 = \pm1$. Jika $q^2 = 1$ maka $(q - 1)(q + 1) = q^2 -
1 = 0$ (sebab skalar pusatnya, sehingga pemfaktoran ini sahih) dan
$q = \pm1$ pada ring pembagi $\mathbb H$, yang memberikan $\rho_q
= I$, jadi terkecualikan; sedangkan $q^2 = -1$ dengan $q = a +
su$ memberikan $a^2 - s^2 + 2as\,u = -1$, sehingga $a = 0$ dan $s = 1$:
jadi $q$ berupa murni satuan $w$, dan $\rho_w$ adalah setengah putaran terhadap
$w$ (bersudut $\pi$, pertanyaan 5). Untuk pusatnya: jika $\rho_q$ komutatif
dengan setiap $\rho_p$, maka $\rho_{pq\bar p} = \rho_q$, sehingga
$pq\bar p = \varepsilon(p)\,q$ dengan $\varepsilon(p) \in
\{\pm1\}$; lalu $p \mapsto \varepsilon(p) = (pq\bar p)q^{-1}$ bersifat
\omterm{def:b3:topology:continuity}{kontinu} pada $S^3$ yang \omterm{def:b3:topology:connected}{terhubung} dan bernilai $1$ di $p =
1$, sehingga $\equiv 1$: jadi $q$ komutatif dengan seluruh $S^3$, sehingga
dengan seluruh $\mathbb H$ (lewat penskalaannya), jadi $q \in \R \cap S^3 =
\{\pm1\}$ (pertanyaan 1) dan $\rho_q = I$: sehingga pusatnya
trivial.

\textbf{22.} Karena $u \perp w$ berupa murni satuan, $uw =
u\wedge w$ bersifat murni (pertanyaan 18), sehingga
\[
w' = qw = \cos\tfrac\theta2\,w + \sin\tfrac\theta2\,u\wedge
w
\]
bersifat murni, bernorma $\abs q\abs w = 1$. Lalu $w'w = qw\cdot w
= qw^2 = -q$, dan
\[
\rho_{w'}\rho_w = \rho_{w'w} = \rho_{-q} = \rho_q .
\]
Sedangkan kedua sumbunya $w$ dan $w' = \cos\frac\theta2 w +
\sin\frac\theta2(u\wedge w)$ terletak pada bidang $u^\perp$
yang ortogonal terhadap sumbu rotasinya, dan $\langle w', w\rangle
= \cos\frac\theta2$: sehingga keduanya bersudut separuh
$\frac\theta2$. Inilah pembangkitan klasiknya: bahwa dua
setengah putaran terhadap sumbu yang bersudut $\frac\theta2$
berkomposisi menjadi rotasi bersudut $\theta$ terhadap tegak lurus
bersamanya.

\textbf{23.} Kekompakannya adalah pertanyaan 7. Untuk keterhubungan \omterm{def:b3:topology:pathconnected}{lintasannya}:
$SO(3) = \rho(S^3)$ merupakan peta \omterm{def:b3:topology:continuity}{kontinu} bola yang
\omterm{def:b3:topology:pathconnected}{terhubung lintasan}; sedangkan sebagai alternatif, untuk $R = \eu^{A}$
dengan $A$ yang antisetangkup (pertanyaan 8), $t \mapsto \eu^{tA}$
merupakan \omterm{def:b3:topology:pathconnected}{lintasan} di $SO(3)$ dari $I$ ke $R$ (yang ortogonal sebab
$(\eu^{tA})^{\mathsf T} = \eu^{-tA}$, dengan determinan $1$ lewat
\omterm{def:b3:topology:continuity}{kekontinuannya} dari $t = 0$). Untuk $O(3)$: $\det$ bersifat \omterm{def:b3:topology:continuity}{kontinu}
ke $\{\pm1\}$, sehingga $O(3)$ tak \omterm{def:b3:topology:connected}{terhubung}, $O(3) = SO(3)
\sqcup D\,SO(3)$ untuk sembarang $D$ yang tetap dengan $\det D = -1$
(misalnya $D = -I$), dan perkalian kiri oleh $D$ merupakan
\omterm{def:b3:topology:continuity}{homeomorfisma}: jadi tepat dua komponen, yang masing-masingnya salinan
$SO(3)$. Jadi $\rho$ yang dua-ke-satu, yang bebas seksi menurut pertanyaan
19, merupakan selimut ganda yang jujur atas sebuah grup \omterm{def:b3:topology:compact}{kompak}
\omterm{def:b3:topology:connected}{terhubung} oleh $S^3$ yang \omterm{def:b3:conformal:simplyconnected}{terhubung sederhana} --- yakni geometri di balik
kiat sabuknya.

\textbf{24.} (i) Untuk matriksnya:
\[
R_1 = \begin{pmatrix} 0 & -1 & 0\\ 1 & 0 & 0\\ 0 & 0 & 1
\end{pmatrix},
\quad
R_2 = \begin{pmatrix} 1 & 0 & 0\\ 0 & 0 & -1\\ 0 & 1 & 0
\end{pmatrix},
\quad
R_2R_1 = \begin{pmatrix} 0 & -1 & 0\\ 0 & 0 & -1\\ 1 & 0 &
0\end{pmatrix}.
\]
Jejaknya $0 = 1 + 2\cos\theta$ memberikan $\cos\theta = -\frac12$:
sehingga $\theta = \frac{2\pi}3$. Sedangkan bagian antisetangkupnya $\frac{R -
R^{\mathsf T}}2$ berentri yang menyandikan sumbunya secara
$\sin\theta\,(v_3, -v_2, v_1)$: di sini $\frac{R - R^{\mathsf T}}2 =
\frac12\bigl(\begin{smallmatrix}0 & -1 & -1\\ 1 & 0 & -1\\ 1
& 1 & 0\end{smallmatrix}\bigr)$, yang terbaca (lewat kamus $A_v$
Bagian V) sebagai $v\sin\theta = \frac12(1, -1, 1)$; lalu dengan
$\sin\frac{2\pi}3 = \frac{\sqrt3}2$: $v = \frac1{\sqrt3}(1,
-1, 1)$.
(ii) Untuk kuaternionnya: $q_1 = \cos\frac\pi4 + \sin\frac\pi4\,k =
\frac{\sqrt2}2(1 + k)$, $q_2 = \frac{\sqrt2}2(1 + i)$, dan
\[
q_2q_1 = \tfrac12(1 + i)(1 + k) = \tfrac12(1 + k + i + ik)
= \tfrac12\bigl(1 + i - j + k\bigr)
\]
(sebab $ik = -j$). Jadi $\cos\frac\theta2 = \frac12$: sehingga $\theta =
\frac{2\pi}3$, dan bagian vektornya $\frac12(i - j + k)$ berarah
$\frac1{\sqrt3}(1, -1, 1)$ --- yakni jawaban yang sama,
dengan jalur kuaternionnya hanya memerlukan satu baris
perkalian alih-alih sebuah hasil kali matriks: yakni alasan praktis
perangkat lunak penerbangan mengomposisikan sikapnya di $S^3$.

\textbf{25.} Untuk keterbalikan $I + K$: $(I + K)v = 0$ memberikan $0 =
\langle v, v\rangle + \langle v, Kv\rangle = \norm v^2$
(sebab keantisetangkupannya membunuh suku keduanya): jadi $v = 0$. Untuk keortogonalan
$C = C(K)$: dengan memakai $(I \pm K)^{\mathsf T} = I \mp K$ dan
fakta bahwa keempat matriks $I \pm K$ dan $(I \pm K)^{-1}$
komutatif (sebab ungkapan polinomial dalam $K$, ditambah limitnya):
\[
C^{\mathsf T}C = (I + K)^{-\mathsf T}(I - K)^{\mathsf T}
(I - K)(I + K)^{-1}
= (I - K)^{-1}(I + K)(I - K)(I + K)^{-1} = I .
\]
Untuk determinannya: $\det(I - K) = \det\bigl((I - K)^{\mathsf
T}\bigr) = \det(I + K)$, sehingga $\det C = 1$: jadi $C \in SO(n)$.
Tak ada nilai eigen $-1$: sebab $Cv = -v$ berarti $(I - K)w = -(I + K)w$
untuk $w = (I + K)^{-1}v$, yakni $2w = 0$: jadi $v = 0$.
Untuk pembalikannya: dari $C(I + K) = I - K$, selesaikanlah $K(I + C) = I -
C$; lalu karena $-1 \notin \operatorname{Sp}C$, $I + C$ bersifat
terbalikkan dan $K = (I - C)(I + C)^{-1}$, yang bersifat
antisetangkup setiap kali $C$ ortogonal tanpa nilai eigen
$-1$ (transposkanlah ungkapannya lalu pakai $C^{\mathsf T} =
C^{-1}$: $K^{\mathsf T} = (I - C^{-1})(I + C^{-1})^{-1} =
(C - I)(C + I)^{-1} = -K$). Jadi kedua pemetaannya saling
berbalikan menurut konstruksinya: yakni sebuah parametrisasi rasional global
atas kepingan terbuka padat $SO(n)$ yang menghindari nilai eigen $-1$
--- tanpa deret, tanpa trigonometri, dan dalam dimensi $3$ ia
merupakan substitusi setengah sudut $K = \tan\frac\theta2\,A_v$ yang menyamar.

\end{solution}
